\documentclass{tsp-short}

\usepackage{amsthm}

\newtheorem{assumption}{Assumption}
\newtheorem{theorem}{Theorem}[section]
\newtheorem{lemma}{Lemma}[section]
\newtheorem{corollary}{Corollary}[section]
\newtheorem{proposition}{Proposition}[section]

\theoremstyle{remark}
\newtheorem{remark}{Remark}[section]

\theoremstyle{definition}
\newtheorem{definition}{Definition}[section]
\newtheorem{example}{Example}[section]

\newcommand{\R}{\mathbb{R}}
\newcommand{\cP}{\mathcal{P}}
\newcommand{\E}{\mathbb E}
\newcommand{\supp}{\mathrm{supp}}
\newcommand{\op}{\mathrm{op}}
\newcommand{\HS}{\mathrm{HS}}
\DeclareMathOperator{\Ent}{Ent}
\DeclareMathOperator{\Var}{Var}
\DeclareMathOperator{\Lip}{Lip}

\usepackage[unicode,hidelinks]{hyperref}
\usepackage[nameinlink,noabbrev]{cleveref} % for cref

\crefname{assumption}{Assumption}{Assumptions}
\Crefname{assumption}{Assumption}{Assumptions}

% \title[Interacting SDEs]{Functional Inequalities for Solutions of Stochastic Differential Equations with Interaction for Non-Constant Diffusion}
\title[Interacting SDEs]{Functional Inequalities and Synchronization for the measure-valued solutions of Stochastic Differential Equations with Interaction}

\author{Kyrylo Kuchynskyi}
\address{Ukraine, Kyiv-4, 3, Tereschenkivska st.}
\email{kkuchynskyi@imath.kiev.ua}
\thanks{This work was supported by a grant from the Simons Foundation (1290607, K.O.K.).}

\subjclass[2020]{60H10, 60B05, 60G57}
\keywords{common noise, interacting SDEs, stochastic flows, logarithmic Sobolev inequality, Jacobian matrix flow}

\begin{document}
% TODO: Rewrite abstract
\begin{abstract}
    We study the measure-valued solution to the stochastic differential equations with interaction.
    we derive pathwise transfer estimates for logarithmic Sobolev, Poincar?, concentration, and \( T_2 \) inequalities.
    The resulting constants are controlled by the maximal Jacobian norm or the Lipschitz distortion of the flow on the
    support of the initial measure. Under global spatial regularity assumptions,
    we obtain moment estimates for the maximal Jacobian norm on compact sets.
    We also establish an averaged Jacobian-energy bound under one-sided derivative assumptions.
    Finally, for a kernel-type interacting model satisfying a pairwise dissipativity condition,
    we prove exponential decay of the expected pairwise variance, yielding mean-square synchronization around
    the random center of mass.
\end{abstract}

\maketitle

% Navigation: generated table of contents for PDF/source orientation.
\tableofcontents


% ============================================================
% Introduction: main SDEWI model and global assumptions
% ============================================================
\section{Introduction}\label{sec:introduction}

% The introduction of SDEWI
We consider stochastic differential equations with interaction that were proposed by Andrey A. Dorogovtsev in \cite{dorogovtsev2003},
%  Main SDEWIN equation
\begin{equation}\label{eq:main-sdewi}
    \begin{cases}
        dX_t(u) & = a\big(X_t(u),\mu_t\big)\,dt + \sigma\big(X_t(u)\big)\,dB_t, \\[1mm]
        X_0(u)  & = u,                                                          \\[1mm]
        \mu_t   & = \mu_0\circ X_t^{-1}.
    \end{cases}
\end{equation}
Here \(\mu_0\) is a deterministic initial probability measure, \(a:\R^d\times\cP_2(\R^d)\to\R^d\) is the drift, where
\(\cP_2(\R^d)\) denotes the space of Borel probability measures on \(\R^d\) with finite second moment.
and the diffusion coefficient \(\sigma:\R^d\to\R^{d\times d}\) is state-dependent and non-constant. Let
\( (\Omega,\mathcal F,(\mathcal F_t)_{t\ge0},\mathbb P) \) be a filtered probability space satisfying the usual
conditions and carrying a \(d\)-dimensional Brownian motion \(B=(B^1,\dots,B^d) \).

% Motivation
The main difference from an ordinary stochastic differential equation is the dependence on all other particles via the push-forward measure \(\mu_t\); that is why
"with the interaction" comes.  The Brownian motion generates the \(\sigma\)-algebra \(\mathcal{F}_t\). For each
realization \(\omega \in \mathcal{F}_t\), the map
\begin{equation*}
    T_t(\omega,\cdot):u\mapsto X_t(u,\omega)
\end{equation*}
acts on the initial measure by pushforward,
\begin{equation*}
    \mu_t(\omega)=\mu_0\circ T_t(\omega)^{-1},
\end{equation*}
so the evolution is best viewed as a random transport of measures by a stochastic flow.

% Main results
We establish three types of estimates for measure-valued solutions generated by the stochastic differential equations
with interactions. First, we prove pathwise transfer results for logarithmic Sobolev, Poincar?, concentration, and
\( T_2\) inequalities under the random transport \( T_t\). Second, under spatial regularity assumptions, we bound
moments of the maximal Jacobian norm on compact sets and derive an averaged Jacobian-energy estimate under
weaker one-sided conditions. Third, for a kernel-type interacting equation satisfying pairwise dissipativity,
we prove exponential decay of the expected pairwise variance and hence mean-square synchronization.

%%% Section 2 motivation %%%
% Matrix norm
In \Cref{sec:jacobian-bound} we investigate the Jacobian norm and how it affects the functional inequalities.
We write \(|\cdot|\) for the Euclidean norm on \(\R^d\).
For a matrix \(M\in\R^{d\times d}\) we set
\begin{equation*}
    \|M\|_{\op}:=\sup_{|v|=1}|Mv|,
    \qquad
    \|M\|_{\HS}:=\Big(\sum_{i,j=1}^d |M_{ij}|^2\Big)^{1/2}.
\end{equation*}

% Assumption on the coefficients
For \(\sigma=(\sigma^{(1)},\ldots,\sigma^{(d)})\) with columns \(\sigma^{(k)}:\R^d\to\R^d\), we define
\begin{equation*}
    \|\partial\sigma(x)\|_{\op}:=\max_{1\le k\le d}\|\partial \sigma^{(k)}(x)\|_{\op}.
\end{equation*}
% Assumption on the coefficients
\begin{assumption}\label{ass:general-lipschitz}
    The coefficients satisfy the following conditions:
    \begin{enumerate}
        \item[(i)] For each \(\mu\in\cP_2(\R^d)\), the map \(x\mapsto a(x,\mu)\) belongs to \(C^2(\R^d;\R^d)\) and
              \begin{equation*}
                  \sup_{x,\mu}\|\partial_x a(x,\mu)\|_{\op}\le L.
              \end{equation*}
        \item[(ii)] The diffusion coefficient satisfies \(\sigma\in C^2(\R^d;\R^{d\times d})\) and there exists \(L_\sigma<\infty\) such that
              \begin{equation*}
                  \sup_x \|\partial \sigma(x)\|_{\op}\le L_\sigma.
              \end{equation*}
        \item[(iii)] There exists \(L_\mu<\infty\) such that for all \(x\in\R^d\) and all \(\mu,\nu\in\cP_2(\R^d)\),
              \begin{equation*}
                  |a(x,\mu)-a(x,\nu)|\le L_\mu\,W_2(\mu,\nu).
              \end{equation*}
    \end{enumerate}
\end{assumption}

% The remarks on the solution
The solution to \cref{eq:main-sdewi} generates a \(C^1\) stochastic flow of maps: there exists a modification such that
for a.s. \(\omega\) and every \(t\ge 0\) the map \(u\mapsto X_t(u,\omega)\) is of class \(C^1(\R^d)\).
We study whether logarithmic Sobolev, Poincar\'e, and Talagrand transport inequalities are propagated from
the initial measure \(\mu_0\) to the random measure \(\mu_t\).


% Logarithmic Sobolev inequality
\begin{definition}[Logarithmic Sobolev inequality]
    Denoting by \(\mathcal{C}^1_c(\R^d)\) the set of continuously differentiable compactly supported functions from \(\R^d\) to \(\R\), for \(g\ge 0\) with \(\int g\,d\mu=1\) we define the entropy
    \begin{equation*}
        \Ent_\mu(g):=\int g\log g\,d\mu.
    \end{equation*}
    A probability measure \(\mu\) on \(\R^d\) is said to satisfy a logarithmic Sobolev inequality with constant \(C>0\) if for every \(h\in\mathcal{C}^1_c(\R^d)\) with \(\int h^2\,d\mu=1\) one has
    \begin{equation}\label{eq:lsi-def}
        \Ent_\mu(h^2)\le C\int_{\R^d}|\nabla h|^2\,d\mu.
    \end{equation}
\end{definition}


% Poincare inequality
\begin{definition}[Poincar\'e inequality]
    A probability measure \(\mu\) on \(\R^d\) is said to satisfy a Poincar\'e inequality with constant \(C_P>0\), denoted by \(\mathrm{PI}(C_P)\), if for every \(h\in \mathcal C_c^1(\R^d)\),
    \begin{equation}\label{eq:poincare-def}
        \Var_\mu(h):=\int_{\R^d} h^2\,d\mu-\Big(\int_{\R^d} h\,d\mu\Big)^2
        \le
        C_P\int_{\R^d} |\nabla h|^2\,d\mu.
    \end{equation}
\end{definition}


% Talagrand transport inequality
\begin{definition}[Talagrand transport inequality]
    Let \(p\ge1\). A probability measure \(\mu\) on \(\R^d\) is said to satisfy the Talagrand transport inequality \(T_p(C_T)\) with constant \(C_T>0\) if for every probability measure \(\nu\) on \(\R^d\) with finite \(p\)-th moment,
    \begin{equation}\label{eq:talagrand-def}
        W_p(\nu,\mu)^2 \le 2C_T\,\Ent(\nu\mid \mu),
    \end{equation}
    where
    \begin{equation*}
        \Ent(\nu\mid\mu)=
        \begin{cases}
            \displaystyle \int_{\R^d}\log\Big(\frac{d\nu}{d\mu}\Big)\,d\nu, & \nu\ll \mu,       \\[3mm]
            +\infty,                                                        & \text{otherwise}.
        \end{cases}
    \end{equation*}
\end{definition}

%%% Section 3 summary %%%
In \Cref{sec:averaged-jacobian}, we replace the pathwise supremum over a compact set by an averaged Jacobian energy and derive a corresponding moment estimate.

%%% Section 4 summary %%%
In \Cref{sec:internal-fluctuations}, we consider a specific form of the SDEWI in which the dependence on \( \mu_t \)
remains explicit through the kernel \( K \) and the parameter \( \Theta_t \).
\begin{equation}\label{eq:kernel-sdewi}
    \begin{cases}
        dX_t(u) & = \bigg[ b(X_t(u),\Theta_t) + \int_{\R^d} K(X_t(u),X_t(v),\Theta_t)\,\mu_0(dv) \bigg]dt \\
                & \quad+ \sum_{\ell=1}^m \sigma_\ell(X_t(u),\Theta_t)\,dB_t^\ell,                         \\[1mm]
        X_0(u)  & =u,\qquad u\in\R^d,                                                                     \\[1mm]
        \mu_t   & =\mu_0\circ X_t^{-1}.
    \end{cases}
\end{equation}
where
\begin{equation*}
    \theta(\mu):=\int_{\R^d}\Phi(x)\,\mu(dx),
\end{equation*}

% Dissipativity Assumptions on global Lipschitz and growth hypotheses.
\begin{assumption}[Global Lipschitz and growth hypotheses]\label{ass:global-lipschitz}
    The maps \(b,K,\sigma_\ell,\Phi\) are globally Lipschitz and have at most linear growth.
    More explicitly, there exists \(L<\infty\) such that
    \begin{align*}
        |\Phi(x)-\Phi(y)|           & \le L|x-y|,                               \\
        |b(x,\theta)-b(y,\eta)|     & \le L\big(|x-y|+|\theta-\eta|\big),       \\
        |K(x,z,\theta)-K(y,w,\eta)| & \le L\big(|x-y|+|z-w|+|\theta-\eta|\big),
    \end{align*}
    and
    \begin{equation*}
        \sum_{\ell=1}^m |\sigma_\ell(x,\theta)-\sigma_\ell(y,\eta)|^2 \le L^2\big(|x-y|^2+|\theta-\eta|^2\big).
    \end{equation*}
\end{assumption}





% ============================================================
% Pathwise Jacobian Bounds and Functional Inequality Transfer
% ============================================================
\section{Pathwise Jacobian Bounds and Functional Inequality Transfer}\label{sec:jacobian-bound}

% Motivation for this section
This section is devoted to the investigation of the SDE for the Jacobian matrix in \cref{eq:jacobian-sde} and the Jacobian norm \(\|J_t(u)\|\) on compact sets.
Such norm is the upper bound for Lipschitz constant of the map \(T_t(u)\). Then knowing the upper bound for Lipschitz constant we can evaluate
the concentration inequality, Poincar\'e inequality, logarithmic Sobolev inequality and Talagrand inequality.


For a fixed \(\omega\) and \(t\), define the flow map \(T_t(\omega,u):=X_t(u,\omega)\) and its Jacobian matrix
\begin{equation*}
    J_t(u):=\partial_u T_t(u)=\partial_u X_t(u)\in\R^{d\times d}.
\end{equation*}
We use the Jacobian SDE from \cite{dorogovtsev2023intermittency}.
For \cref{eq:main-sdewi}, it has the following form.

\begin{theorem}[\cite{dorogovtsev2023intermittency}]\label{thm:jacobian-sde}
    Under \Cref{ass:general-lipschitz}, for each \(u\in\R^d\), the Jacobian matrix process \(J_t(u)\) exists and satisfies the linear matrix SDE (cf.\ \cite{Kunita1986,dorogovtsev2023intermittency})
    \begin{equation}\label{eq:jacobian-sde}
        dJ_t(u)= A_t(u)\,J_t(u)\,dt + \sum_{k=1}^d C_t^{(k)}(u)\,J_t(u)\,dB_t^{k},\qquad J_0(u)=I,
    \end{equation}
    where
    \begin{equation*}
        A_t(u):=\partial_x a\big(X_t(u),\mu_t\big),\qquad
        C_t^{(k)}(u):=\partial_x \sigma^{(k)}\big(X_t(u)\big),
    \end{equation*}
    and \(\sigma^{(k)}\) denotes the \(k\)-th column of \(\sigma\).
    Here \(A_t(u),C_t^{(k)}(u)\in\R^{d\times d}\), the driving Brownian motion is \(B_t=(B_t^1,\ldots,B_t^d)\) with \(B_t^k\in\R\), and the sum runs over the \(d\) components of the noise.
    Moreover, the coefficients satisfy the uniform bounds
    \begin{equation*}
        \|A_t(u)\|_{\op}\le L,\qquad
        \|C_t^{(k)}(u)\|_{\op}\le L_\sigma
        \quad \text{for all }t,u,\omega.
    \end{equation*}
\end{theorem}

\begin{corollary}[Moment bound for bounded linear SDEs]\label{cor:moment-bound}
    Let \(Y_t\) solve
    \begin{equation*}
        dY_t=A_tY_t\,dt+\sum_{k=1}^d C_t^{(k)}Y_t\,dB_t^k,
        \qquad Y_0=y,
    \end{equation*}
    where \(\|A_t\|_{\op}\le L\) and \(\|C_t^{(k)}\|_{\op}\le L_\sigma\). Then, for every \(q\ge1\),
    \begin{equation*}
        \E|Y_t|^{2q} \le |y|^{2q}\exp\Big(q(2L+dL_\sigma^2)t+2q^2dL_\sigma^2t\Big).
    \end{equation*}
    Consequently, the Jacobian flow satisfies
    \begin{equation*}
        \E\,\|J_t(u)\|_{\op}^{2q} \le d^q \exp\Big(q(2L+dL_\sigma^2)t+2q^2dL_\sigma^2t\Big).
    \end{equation*}
\end{corollary}

\noindent\emph{Proof.}
Apply It\^o's formula to \(|Y_t|^2\), use the bounds on \(A_t\) and
\(C_t^{(k)}\), and then apply the exponential martingale estimate. The matrix
estimate follows by applying the vector estimate to the columns of \(J_t(u)\)
and using \(\|J_t(u)\|_{\op}\le \|J_t(u)\|_{\HS}\).
\hfill\(\Box\)

This theorem shows that even the drift coefficient depends on the measure \( \mu_t \), the Jacobian SDE in \cref{eq:jacobian-sde} does not contain the
derivative with respect to it. Because of that fact we can evaluate \(\|J_t(u)\|\) using Kunita's theorem for the stochastic flows.

Under \Cref{ass:general-lipschitz}, standard results on stochastic flows (see, e.g., \cite[Theorem~1.4.1]{Kunita1986}) yield a modification such that for a.s.\ \(\omega\) and each \(t\ge0\)
the map \(u\mapsto T_t(\omega,u)\) is \(\mathcal{C}^1\).



\begin{definition}[Random Jacobian constant]\label{def:jacobian-constant}
    For a nonempty set \(K\subset\R^d\), define
    \begin{equation*}
        \mathcal{K}_t(\omega;K):=\sup_{u\in K}\|J_t(u,\omega)\|_{\op}.
    \end{equation*}
    If \(K\) is compact and \(u\mapsto J_t(u,\omega)\) is continuous, then \(\mathcal{K}_t(\omega;K)<\infty\).
\end{definition}

\begin{remark}[Why \(\mathcal{K}_t(\omega;K)\) is finite and relation to global blow-up examples]
    The counterexamples of Mohammed--Scheutzow show that even for globally Lipschitz coefficients one may fail to have a finite
    global constant \(\sup_{u\in\R^d}\|J_t(u,\omega)\|_{\op}\) for a fixed \((t,\omega)\).
    In contrast, the quantity \(\mathcal{K}_t(\omega;K)\) is taken over a \emph{compact} set \(K\); under \Cref{ass:general-lipschitz}(iii),
    for each fixed \((t,\omega)\) the map \(u\mapsto J_t(u,\omega)\) is continuous, hence the supremum over \(K\) is finite.
\end{remark}

\begin{corollary}[Pathwise LSI with the Jacobian matrix constant]\label{cor:pathwise-lsi}
    If \(\mu_0\) satisfies \(\mathrm{LSI}(C_0)\) and \(\supp(\mu_0)\subset K\) for some compact \(K\), then for every \(t\ge0\) and a.s.\ \(\omega\),
    \begin{equation*}
        \mu_t(\omega)=\mu_0\circ T_t(\omega)^{-1}
        \quad\Longrightarrow\quad
        \mu_t(\omega)\text{ satisfies }\mathrm{LSI}\big(C_0\,\mathcal{K}_t(\omega;K)^2\big).
    \end{equation*}
\end{corollary}

\begin{lemma}[Flow Lipschitz constant via Jacobian matrix]\label{lem:lip-via-jacobian}
    For each fixed \((t,\omega)\) such that the flow map \(u\mapsto T_t(\omega,u):=X_t(u,\omega)\) is of class \(C^1\),
    \begin{equation*}
        \Lip(T_t(\omega)) = \sup_{u\in\R^d}\|J_t(u,\omega)\|_{\op}\in[0,\infty].
    \end{equation*}
\end{lemma}

\noindent\emph{Proof.}
Since \(u\mapsto T_t(\omega,u)\) is \(C^1\), the mean value theorem implies
\begin{equation*}
    |T_t(u)-T_t(v)| \le \sup_{\theta\in[0,1]}\|J_t(v+\theta(u-v))\|_{\op}\,|u-v|.
\end{equation*}
Taking the supremum over \(u\neq v\) yields \(\Lip(T_t)\le \sup_u\|J_t(u)\|_{\op}\).
Conversely, for any \(u\) and unit vector \(\xi\), letting \(v=u+\varepsilon\xi\) and sending \(\varepsilon\downarrow 0\) gives
\begin{equation*}
    \|J_t(u)\|_{\op} = \sup_{|\xi|=1}\lim_{\varepsilon\downarrow 0}\frac{|T_t(u+\varepsilon\xi)-T_t(u)|}{\varepsilon}
    \le \Lip(T_t),
\end{equation*}
hence equality.
\hfill\(\Box\)




Pointwise estimates control \(\|J_t(u)\|_{\op}\) for fixed \(u\). To control the random Lipschitz constant
\begin{equation*}
    \mathcal{K}_t(\omega;K):=\sup_{u\in K}\|J_t(u,\omega)\|_{\op}
\end{equation*}
on a compact set \(K\), we impose the following additional spatial regularity assumption.
\begin{assumption}\label{ass:spatial-lip-derivatives}
    There exist constants \(L^{(2)}<\infty\) and \(L_\sigma^{(2)}<\infty\) such that for all \(\mu\in\cP_2(\R^d)\) and all \(x,y\in\R^d\),
    \begin{align*}
        \|\partial_x a(x,\mu)-\partial_x a(y,\mu)\|_{\op}                             & \le L^{(2)}\,|x-y|,        \\
        \max_{1\le k\le d}\|\partial \sigma^{(k)}(x)-\partial \sigma^{(k)}(y)\|_{\op} & \le L_\sigma^{(2)}\,|x-y|.
    \end{align*}
\end{assumption}

\begin{proposition}[Moment bound for \(\mathcal{K}_t(\cdot;K)\)]\label{prop:kt-moment}
    Assume \Cref{ass:general-lipschitz,ass:spatial-lip-derivatives}. Then for every compact \(K\subset\R^d\), every \(p\ge 1\), and every \(t\ge0\),
    \begin{equation}\label{eq:kt-finite-moment}
        \E\,\mathcal{K}_t(\cdot;K)^p<\infty.
    \end{equation}
    Moreover, there exist constants \(C_{p,K},c_p<\infty\) (depending only on \(p,K,d,L,L_\sigma,L^{(2)},L_\sigma^{(2)}\)) such that
    \begin{equation}\label{eq:kt-exp-growth}
        \E\,\mathcal{K}_t(\cdot;K)^p \le C_{p,K}\,e^{c_p t}, \qquad t\ge0.
    \end{equation}
\end{proposition}



\noindent\emph{Proof.}
Fix a compact set $K\subset\R^d$, $t\ge0$, and $p\ge1$. Choose an even integer $m>\max\{2p,2d\}$. It is enough to prove \eqref{eq:kt-exp-growth} with $p$ replaced by $m$, because then H\"older's inequality gives
\begin{equation*}
    \E\,\mathcal{K}_t(\cdot;K)^p \le \big(\E\,\mathcal{K}_t(\cdot;K)^m\big)^{p/m}.
\end{equation*}
Fix \(u,v\in\R^d\) and set
\begin{equation*}
    \Delta X_t:=X_t(u)-X_t(v).
\end{equation*}
Since the same random measure \(\mu_t\) appears in both equations, subtraction yields
\begin{equation*}
    d\Delta X_t = \big(a(X_t(u),\mu_t)-a(X_t(v),\mu_t)\big)\,dt + \big(\sigma(X_t(u))-\sigma(X_t(v))\big)\,dB_t.
\end{equation*}
For \(u\neq v\), define matrix-valued processes
\begin{equation*}
    \widetilde A_t(u,v)
    :=
    \int_0^1
    \partial_x a\big(X_t(v)+\theta(X_t(u)-X_t(v)),\mu_t\big)\,d\theta
\end{equation*}
and, for \(k=1,\dots,d\),
\begin{equation*}
    \widetilde C_t^{(k)}(u,v) := \int_0^1  \partial_x \sigma^{(k)}\big(X_t(v)+\theta(X_t(u)-X_t(v))\big)\,d\theta,
\end{equation*}
where \(\sigma^{(k)}\) denotes the \(k\)-th column of \(\sigma\).
By the fundamental theorem of calculus,
\begin{equation*}
    a(X_t(u),\mu_t)-a(X_t(v),\mu_t) = \widetilde A_t(u,v)\,\Delta X_t,
\end{equation*}
and
\begin{equation*}
    \sigma^{(k)}(X_t(u))-\sigma^{(k)}(X_t(v))
    =
    \widetilde C_t^{(k)}(u,v)\,\Delta X_t.
\end{equation*}
Hence \(\Delta X_t\) satisfies the linear SDE
\begin{equation*}
    d\Delta X_t
    =
    \widetilde A_t(u,v)\,\Delta X_t\,dt
    +
    \sum_{k=1}^d \widetilde C_t^{(k)}(u,v)\,\Delta X_t\,dB_t^k,
    \qquad
    \Delta X_0=u-v.
\end{equation*}
Moreover, by Assumption~\ref{ass:general-lipschitz},
\begin{equation*}
    \|\widetilde A_t(u,v)\|_{\op}\le L,
    \qquad
    \|\widetilde C_t^{(k)}(u,v)\|_{\op}\le L_\sigma.
\end{equation*}
If \(u=v\), then \(\Delta X_t\equiv0\), so the desired estimate is trivial. Assume therefore that \(u\neq v\), and set
\begin{equation*}
    Z_t:=\frac{\Delta X_t}{|u-v|},
    \qquad
    \xi:=\frac{u-v}{|u-v|}.
\end{equation*}
Then \(|\xi|=1\), \(Z_0=\xi\), and \(Z_t\) solves
\begin{equation*}
    dZ_t
    =
    \widetilde A_t(u,v)\,Z_t\,dt
    +
    \sum_{k=1}^d \widetilde C_t^{(k)}(u,v)\,Z_t\,dB_t^k.
\end{equation*}
Applying Corollary~\ref{cor:moment-bound} to this linear equation with exponent \(q=m/2\), we obtain
\begin{equation*}
    \E|Z_t|^m
    \le
    C_m e^{c_m t},
\end{equation*}
for constants \(C_m,c_m<\infty\) depending only on \(m,d,L,L_\sigma\).
Multiplying by \(|u-v|^m\) yields
\begin{equation}\label{eq:flow-increment-moment}
    \E\,|X_t(u)-X_t(v)|^m
    \le
    C_m e^{c_m t}|u-v|^m,
    \qquad u,v\in\R^d.
\end{equation}




By Corollary~\ref{cor:moment-bound}, for every $q\ge1$ and every $u\in\R^d$,
\begin{equation*}
    \E\,\|J_t(u)\|_{\op}^{2q}
    \le d^q \exp\Big(q(2L+dL_\sigma^2)t+2q^2 dL_\sigma^2 t\Big).
\end{equation*}
In particular, taking $q=m/2$ and $q=m$, we obtain
\begin{equation}\label{eq:jacobian-moment}
    \E\,\|J_t(u)\|_{\op}^m
    \le C_m e^{c_m t},
    \qquad
    \E\,\|J_t(u)\|_{\op}^{2m}
    \le C_m e^{c_m t},
\end{equation}
uniformly in $u\in\R^d$.




Fix \(u,v\in\R^d\) and define
\begin{equation*}
    D_t:=J_t(u)-J_t(v).
\end{equation*}
Subtracting the Jacobian equations \eqref{eq:jacobian-sde} at \(u\) and \(v\), we obtain
\begin{equation}\label{eq:jacobian-difference-sde}
    \begin{aligned}
        dD_t
         & =
        A_t(u)D_t\,dt+\sum_{k=1}^d C_t^{(k)}(u)D_t\,dB_t^k \\
         & \quad
        +\big(A_t(u)-A_t(v)\big)J_t(v)\,dt
        +\sum_{k=1}^d \big(C_t^{(k)}(u)-C_t^{(k)}(v)\big)J_t(v)\,dB_t^k,
        \qquad D_0=0,
    \end{aligned}
\end{equation}
where
\begin{equation*}
    A_t(w):=\partial_x a(X_t(w),\mu_t),
    \qquad
    C_t^{(k)}(w):=\partial_x \sigma^{(k)}(X_t(w)).
\end{equation*}
By Assumptions~\ref{ass:general-lipschitz} and \ref{ass:spatial-lip-derivatives},
\begin{equation*}
    \|A_t(w)\|_{\op}\le L,
    \qquad
    \|C_t^{(k)}(w)\|_{\op}\le L_\sigma,
\end{equation*}
and
\begin{equation}\label{eq:coefficient-difference}
    \|A_t(u)-A_t(v)\|_{\op}
    \le L^{(2)}|X_t(u)-X_t(v)|,
    \qquad
    \|C_t^{(k)}(u)-C_t^{(k)}(v)\|_{\op}
    \le L_\sigma^{(2)}|X_t(u)-X_t(v)|.
\end{equation}




We estimate
\begin{equation*}
    \Phi(t):=\E\sup_{s\le t}\|D_s\|_{\HS}^m.
\end{equation*}
Set
\begin{equation*}
    F_s:=\big(A_s(u)-A_s(v)\big)J_s(v),
    \qquad
    G_s^{(k)}:=\big(C_s^{(k)}(u)-C_s^{(k)}(v)\big)J_s(v).
\end{equation*}
Then
\begin{equation*}
    dD_t
    =
    A_t(u)D_t\,dt+\sum_{k=1}^d C_t^{(k)}(u)D_t\,dB_t^k
    +
    F_t\,dt+\sum_{k=1}^d G_t^{(k)}\,dB_t^k,
    \qquad D_0=0.
\end{equation*}
Hence, for \(s\le t\),
\begin{equation*}
    D_s
    =
    \int_0^s A_r(u)D_r\,dr
    +\sum_{k=1}^d \int_0^s C_r^{(k)}(u)D_r\,dB_r^k
    +\int_0^s F_r\,dr
    +\sum_{k=1}^d \int_0^s G_r^{(k)}\,dB_r^k.
\end{equation*}
Using \((a_1+\cdots+a_4)^m\le C_m(a_1^m+\cdots+a_4^m)\), we obtain
\begin{align*}
    \sup_{s\le t}\|D_s\|_{\HS}^m
     & \le
    C_m\Bigg[
    \sup_{s\le t}\Big\|\int_0^s A_r(u)D_r\,dr\Big\|_{\HS}^m
    +
    \sup_{s\le t}\Big\|\sum_{k=1}^d \int_0^s C_r^{(k)}(u)D_r\,dB_r^k\Big\|_{\HS}^m \\
     & \hspace{2cm}
    +
    \sup_{s\le t}\Big\|\int_0^s F_r\,dr\Big\|_{\HS}^m
    +
    \sup_{s\le t}\Big\|\sum_{k=1}^d \int_0^s G_r^{(k)}\,dB_r^k\Big\|_{\HS}^m
    \Bigg].
\end{align*}
Taking expectations, we estimate the four terms separately.

For the drift term involving $A_r(u)D_r$,
\begin{equation*}
    \sup_{s\le t}\Big\|\int_0^s A_r(u)D_r\,dr\Big\|_{\HS}
    \le
    L\int_0^t \|D_r\|_{\HS}\,dr,
\end{equation*}
hence by Jensen's inequality,
\begin{equation*}
    \E\sup_{s\le t}\Big\|\int_0^s A_r(u)D_r\,dr\Big\|_{\HS}^m
    \le
    L^m t^{m-1}\int_0^t \E\sup_{\rho\le r}\|D_\rho\|_{\HS}^m\,dr.
\end{equation*}
For the stochastic term involving \(C_r^{(k)}(u)D_r\), the BDG inequality yields
\begin{align*}
    \E\sup_{s\le t}\Big\|\sum_{k=1}^d \int_0^s C_r^{(k)}(u)D_r\,dB_r^k\Big\|_{\HS}^m
     & \le
    C_m\,
    \E\Big(\int_0^t \sum_{k=1}^d \|C_r^{(k)}(u)D_r\|_{\HS}^2\,dr\Big)^{m/2} \\
     & \le
    C_m\,
    \E\Big(\int_0^t dL_\sigma^2\,\|D_r\|_{\HS}^2\,dr\Big)^{m/2}.
\end{align*}
Since \(m\ge2\), Jensen's inequality gives
\begin{equation*}
    \Big(\int_0^t \|D_r\|_{\HS}^2\,dr\Big)^{m/2}
    \le
    t^{\frac m2-1}\int_0^t \sup_{\rho\le r}\|D_\rho\|_{\HS}^m\,dr,
\end{equation*}
and therefore
\begin{equation*}
    \E\sup_{s\le t}\Big\|\sum_{k=1}^d \int_0^s C_r^{(k)}(u)D_r\,dB_r^k\Big\|_{\HS}^m
    \le
    C_m t^{\frac m2-1}\int_0^t \Phi(r)\,dr.
\end{equation*}
Next,
\begin{equation*}
    \sup_{s\le t}\Big\|\int_0^s F_r\,dr\Big\|_{\HS}^m
    \le
    t^{m-1}\int_0^t \|F_r\|_{\HS}^m\,dr,
\end{equation*}
hence
\begin{equation*}
    \E\sup_{s\le t}\Big\|\int_0^s F_r\,dr\Big\|_{\HS}^m
    \le
    t^{m-1}\int_0^t \E\|F_r\|_{\HS}^m\,dr.
\end{equation*}
Finally, applying the BDG inequality to the martingale term involving \(G_r^{(k)}\), we obtain
\begin{align*}
    \E\sup_{s\le t}\Big\|\sum_{k=1}^d \int_0^s G_r^{(k)}\,dB_r^k\Big\|_{\HS}^m
     & \le
    C_m\,
    \E\Big(\int_0^t \sum_{k=1}^d \|G_r^{(k)}\|_{\HS}^2\,dr\Big)^{m/2} \\
     & \le
    C_m t^{\frac m2-1}\int_0^t \sum_{k=1}^d \E\|G_r^{(k)}\|_{\HS}^m\,dr.
\end{align*}

Combining the four estimates, we obtain
\begin{align}
    \Phi(t)
     & \le
    C_m(1+t^{m-1}+t^{\frac m2-1})\int_0^t \Phi(r)\,dr \notag \\
     & \quad+
    C_m(1+t^{m-1}+t^{\frac m2-1})\int_0^t
    \Big(
    \E\|F_r\|_{\HS}^m
    +
    \sum_{k=1}^d \E\|G_r^{(k)}\|_{\HS}^m
    \Big)\,dr.
    \label{eq:jacobian-gronwall}
\end{align}

We now estimate the forcing terms. By \eqref{eq:coefficient-difference},
\begin{equation*}
    \|F_r\|_{\HS}
    \le
    L^{(2)}|X_r(u)-X_r(v)|\,\|J_r(v)\|_{\HS},
\end{equation*}
and
\begin{equation*}
    \|G_r^{(k)}\|_{\HS}
    \le
    L_\sigma^{(2)}|X_r(u)-X_r(v)|\,\|J_r(v)\|_{\HS}.
\end{equation*}
Hence, by H\"older's inequality,
\begin{align*}
    \E\|F_r\|_{\HS}^m
     & \le
    C\,
    \Big(\E|X_r(u)-X_r(v)|^{2m}\Big)^{1/2}
    \Big(\E\|J_r(v)\|_{\HS}^{2m}\Big)^{1/2}, \\
    \E\|G_r^{(k)}\|_{\HS}^m
     & \le
    C\,
    \Big(\E|X_r(u)-X_r(v)|^{2m}\Big)^{1/2}
    \Big(\E\|J_r(v)\|_{\HS}^{2m}\Big)^{1/2}.
\end{align*}
Since \(\|M\|_{\HS}\le \sqrt d\,\|M\|_{\op}\), \eqref{eq:flow-increment-moment} and \eqref{eq:jacobian-moment} yield
\begin{equation*}
    \E\|F_r\|_{\HS}^m
    +
    \sum_{k=1}^d \E\|G_r^{(k)}\|_{\HS}^m
    \le
    C_m e^{c_m r}|u-v|^m.
\end{equation*}
Substituting this into \eqref{eq:jacobian-gronwall}, we obtain
\begin{equation*}
    \Phi(t)
    \le
    C_m(1+t^{m-1}+t^{\frac m2-1})\int_0^t \Phi(r)\,dr
    +
    C_m(1+t^{m-1}+t^{\frac m2-1})e^{c_m t}|u-v|^m.
\end{equation*}
By Gronwall's lemma,
\begin{equation*}
    \Phi(t)
    \le
    C_m(1+t^{m-1}+t^{\frac m2-1})e^{c_m t}|u-v|^m.
\end{equation*}
Absorbing the polynomial factor into the exponential by enlarging \(c_m\), we get
\begin{equation*}
    \Phi(t)\le C_m e^{c_m t}|u-v|^m.
\end{equation*}
In particular,
\begin{equation*}
    \E\,\|D_t\|_{\HS}^m
    \le
    \E\sup_{s\le t}\|D_s\|_{\HS}^m
    \le
    C_m e^{c_m t}|u-v|^m.
\end{equation*}
Since \(\|D_t\|_{\op}\le \|D_t\|_{\HS}\), we conclude that
\begin{equation}\label{eq:jacobian-increment}
    \E\,\|J_t(u)-J_t(v)\|_{\op}^m
    \le
    C_m e^{c_m t}|u-v|^m,
    \qquad u,v\in\R^d.
\end{equation}





Fix \(\beta\in(0,1-d/m)\).
Since \(m>d\), the increment estimate \eqref{eq:jacobian-increment} implies, by the Kolmogorov continuity theorem applied in the spatial variable \(u\), that for each fixed \(t\ge0\) there exists a modification of \(u\mapsto J_t(u)\) whose sample paths on \(K\) are \(\beta\)-H\"older continuous.
Moreover, by a quantitative version of the Kolmogorov theorem, one may choose such a modification, still denoted by \(J_t\), so that
\begin{equation}\label{eq:holder-bound}
    \E\,[J_t]_{\beta;K}^m
    \le
    C_{m,\beta,K}e^{c_m t},
\end{equation}
where
\begin{equation*}
    [J_t]_{\beta;K}
    :=
    \sup_{\substack{u,v\in K\\u\neq v}}
    \frac{\|J_t(u)-J_t(v)\|_{\op}}{|u-v|^\beta}.
\end{equation*}
Choose any point \(u_0\in K\).
Then for every \(u\in K\),
\begin{equation*}
    \|J_t(u)\|_{\op}
    \le
    \|J_t(u_0)\|_{\op}
    +
    [J_t]_{\beta;K}|u-u_0|^\beta
    \le
    \|J_t(u_0)\|_{\op}
    +
    [J_t]_{\beta;K}(\operatorname{diam}K)^\beta.
\end{equation*}
Therefore,
\begin{equation*}
    \mathcal{K}_t(\omega;K)
    =
    \sup_{u\in K}\|J_t(u,\omega)\|_{\op}
    \le
    \|J_t(u_0,\omega)\|_{\op}
    +
    [J_t(\omega)]_{\beta;K}(\operatorname{diam}K)^\beta.
\end{equation*}
Taking \(m\)-th moments and using \eqref{eq:jacobian-moment} and \eqref{eq:holder-bound}, we obtain
\begin{equation*}
    \E\,\mathcal{K}_t(\cdot;K)^m
    \le
    C_{m,K}e^{c_m t}.
\end{equation*}
As explained at the beginning, this implies
\begin{equation*}
    \E\,\mathcal{K}_t(\cdot;K)^p
    \le
    C_{p,K}e^{c_p t},
    \qquad t\ge0,
\end{equation*}
for every \(p\ge1\).
In particular, \eqref{eq:kt-finite-moment} holds.
\hfill\(\Box\)






% ============================================================
% Concentration inequality
% ============================================================
\subsection{Concentration inequality}

In this subsection we record several consequences of the identity
\begin{equation*}
    T_t(\omega,u):=X_t(u,\omega),
\end{equation*}
which expresses the random measure \(\mu_t\) as the pushforward of the initial law \(\mu_0\) under the stochastic flow.

We next record the basic concentration-transfer principle under a random Lipschitz map.

\begin{proposition}[Pathwise concentration transfer]\label{prop:concentration-transfer}
    Let \(K:=\supp(\mu_0)\) and assume that for a.s.\ \(\omega\) the map \(T_t(\omega,\cdot)\) is Lipschitz on \(K\), with Lipschitz constant
    \begin{equation*}
        L_t(\omega;K):=\sup_{\substack{u,v\in K\\u\neq v}}
        \frac{|T_t(\omega,u)-T_t(\omega,v)|}{|u-v|}.
    \end{equation*}
    Assume that \(\mu_0\) satisfies the concentration inequality
    \begin{equation}\label{eq:general-concentration-mu0}
        \mu_0(f-\mu_0(f)\ge r)\le \Psi(r),
        \qquad r\ge0,
    \end{equation}
    for every \(1\)-Lipschitz function \(f:\R^d\to\R\), where \(\Psi:[0,\infty)\to[0,1]\) is nonincreasing.
    Then for a.s.\ \(\omega\),
    \begin{equation}\label{eq:general-concentration-mut}
        \mu_t(\omega)(g-\mu_t(\omega)(g)\ge r)
        \le
        \Psi\Big(\frac{r}{L_t(\omega;K)}\Big),
        \qquad r\ge0,
    \end{equation}
    for every \(1\)-Lipschitz function \(g:\R^d\to\R\).
\end{proposition}

\noindent\emph{Proof.}
Fix \(\omega\) such that \(T_t(\omega,\cdot)\) is Lipschitz on \(K\), and write
\(L:=L_t(\omega;K)\in[0,\infty)\).
Let \(g:\R^d\to\R\) be \(1\)-Lipschitz.

If \(L=0\), then \(T_t(\omega,\cdot)\) is constant on \(K\), hence
\(\mu_t(\omega)=\mu_0\circ T_t(\omega)^{-1}\) is a Dirac mass and
\(\mu_t(\omega)(g-\mu_t(\omega)(g)\ge r)=0\) for every \(r>0\).
Thus \cref{eq:general-concentration-mut} is immediate.

Assume now \(L>0\) and define on \(K\)
\begin{equation*}
    h(u):=\frac{g(T_t(\omega,u))}{L}.
\end{equation*}
Then \(h\) is \(1\)-Lipschitz on \(K\).
By the McShane extension theorem, \(h\) extends to a \(1\)-Lipschitz function
\(\widetilde h:\R^d\to\R\).
Applying \cref{eq:general-concentration-mu0} to \(\widetilde h\) (and using that \(\mu_0\) is supported on \(K\)) gives
\begin{equation*}
    \mu_0\big(h-\mu_0(h)\ge s\big)\le \Psi(s),\qquad s\ge0.
\end{equation*}
Since \(\mu_t(\omega)=\mu_0\circ T_t(\omega)^{-1}\),
\begin{equation*}
    \mu_t(\omega)(g)=\mu_0(g\circ T_t(\omega,\cdot))=L\mu_0(h).
\end{equation*}
Therefore, for \(r\ge0\),
\begin{align*}
    \mu_t(\omega)(g-\mu_t(\omega)(g)\ge r)
     & = \mu_0\big(g\circ T_t(\omega,\cdot)-\mu_0(g\circ T_t(\omega,\cdot))\ge r\big) \\
     & = \mu_0\big(h-\mu_0(h)\ge r/L\big)                                             \\
     & \le \Psi\Big(\frac{r}{L}\Big)
    = \Psi\Big(\frac{r}{L_t(\omega;K)}\Big).
\end{align*}
This proves \cref{eq:general-concentration-mut}.
\hfill\(\Box\)

\begin{remark}[Concentration as a corollary of LSI]
    If \(\mu_0\) satisfies a logarithmic Sobolev inequality, then it satisfies Gaussian concentration.
    Hence \Cref{cor:pathwise-lsi} yields, in particular, a pathwise concentration inequality for \(\mu_t(\omega)\) with a random constant controlled by \(\mathcal K_t(\omega;K)^2\).
\end{remark}






% ============================================================
% Poincar\'e inequality
% ============================================================
\subsection{Poincar\'e inequality}\label{sec:poincare}

We first record the stability of the Poincar\'e inequality under Lipschitz
pushforward.

\begin{proposition}[Poincar\'e inequality under a Lipschitz pushforward]\label{prop:poincare-pushforward}
    Let \(\mu\) be a probability measure on \(\R^d\) satisfying \(\mathrm{PI}(C_P)\), and let
    \(T:\R^d\to\R^d\) be a Lipschitz map with Lipschitz constant \(\Lip(T)\).
    Then the pushforward measure \(\nu:=\mu\circ T^{-1}\) satisfies
    \begin{equation*}
        \mathrm{PI}\big(C_P\,\Lip(T)^2\big).
    \end{equation*}
\end{proposition}

\noindent\emph{Proof.}
Let \(g\in \mathcal C_c^1(\R^d)\) and set \(f:=g\circ T\).
Then
\begin{equation*}
    \int_{\R^d} g\,d\nu=\int_{\R^d} g(T(u))\,\mu(du)=\int_{\R^d} f(u)\,\mu(du),
\end{equation*}
and similarly
\begin{equation*}
    \int_{\R^d} g^2\,d\nu=\int_{\R^d} f^2\,d\mu.
\end{equation*}
Hence
\begin{equation*}
    \Var_\nu(g)=\Var_\mu(f).
\end{equation*}
Applying the Poincar\'e inequality for \(\mu\) gives
\begin{equation*}
    \Var_\nu(g) = \Var_\mu(f) \le C_P\int_{\R^d} |\nabla f(u)|^2\,\mu(du).
\end{equation*}
Since \(f=g\circ T\) and \(T\) is Lipschitz, we have a.e.
\begin{equation*}
    |\nabla f(u)| \le \Lip(T)\,|\nabla g(T(u))|.
\end{equation*}
Therefore
\begin{equation*}
    \Var_\nu(g) \le C_P\,\Lip(T)^2 \int_{\R^d} |\nabla g(T(u))|^2\,\mu(du) = C_P\,\Lip(T)^2 \int_{\R^d} |\nabla g(x)|^2\,\nu(dx).
\end{equation*}
This proves the claim.
\hfill\(\Box\)

\begin{corollary}[Pathwise Poincar\'e inequality]\label{cor:pathwise-poincare}
    If \(\mu_0\) satisfies \(\mathrm{PI}(C_P)\) and \(\supp(\mu_0)\subset K\) for some compact set \(K\subset \R^d\), then for every \(t\ge0\) and a.s.\ \(\omega\),
    \begin{equation*}
        \mu_t(\omega)=\mu_0\circ T_t(\omega)^{-1}
        \quad\Longrightarrow\quad
        \mu_t(\omega)\text{ satisfies }\mathrm{PI}\big(C_P\,\mathcal K_t(\omega;K)^2\big).
    \end{equation*}
\end{corollary}

\noindent\emph{Proof.}
By \Cref{prop:poincare-pushforward}, it suffices to bound the Lipschitz constant of \(T_t(\omega,\cdot)\) on \(K\).
Since \(K\) is compact and \(u\mapsto T_t(\omega,u)\) is of class \(C^1\), the mean value theorem gives
\begin{equation*}
    |T_t(\omega,u)-T_t(\omega,v)|
    \le
    \sup_{z\in K}\|J_t(z,\omega)\|_{\op}\,|u-v|
    =
    \mathcal K_t(\omega;K)\,|u-v|,
    \qquad u,v\in K.
\end{equation*}
Hence \(\Lip(T_t(\omega)|_K)\le \mathcal K_t(\omega;K)\), and the result follows.
\hfill\(\Box\)






% ============================================================
% Talagrand transport inequalities
% ============================================================
\subsection{Talagrand transport inequalities}\label{sec:talagrand}

We next record the analogous stability property for Talagrand transport
inequalities.

\begin{proposition}[Transport inequality under a Lipschitz pushforward]\label{prop:talagrand-pushforward}
    Let \(\mu\) be a probability measure on \(\R^d\) satisfying \(T_p(C_T)\) for some \(p\ge1\), and let
    \(T:\R^d\to\R^d\) be a Lipschitz map with Lipschitz constant \(\Lip(T)\).
    Then the pushforward measure \(\eta:=\mu\circ T^{-1}\) satisfies
    \begin{equation*}
        T_p\big(C_T\,\Lip(T)^2\big).
    \end{equation*}
\end{proposition}

\noindent\emph{Proof.}
Let \(\lambda\) be any probability measure on \(\R^d\).
If \(\lambda\not\ll \eta\), then \(\Ent(\lambda\mid\eta)=+\infty\) and the inequality is trivial.
Assume therefore that \(\lambda\ll \eta\).

Choose a probability measure \(\widetilde\lambda\) on \(\R^d\) such that
\begin{equation*}
    \widetilde\lambda\circ T^{-1}=\lambda
    \qquad\text{and}\qquad
    \Ent(\widetilde\lambda\mid \mu)=\Ent(\lambda\mid \eta).
\end{equation*}
For instance, one may take the measure
\begin{equation*}
    \widetilde\lambda(du):=
    \frac{d\lambda}{d\eta}(T(u))\,\mu(du).
\end{equation*}
Indeed,
\begin{equation*}
    (\widetilde\lambda\circ T^{-1})(A)
    =
    \int_{\R^d}\mathbf 1_{\{T(u)\in A\}}\frac{d\lambda}{d\eta}(T(u))\,\mu(du)
    =
    \int_A \frac{d\lambda}{d\eta}(x)\,\eta(dx)
    =
    \lambda(A),
\end{equation*}
and
\begin{equation*}
    \Ent(\widetilde\lambda\mid\mu)
    =
    \int_{\R^d}\log\Big(\frac{d\lambda}{d\eta}(T(u))\Big)\,\widetilde\lambda(du)
    =
    \int_{\R^d}\log\Big(\frac{d\lambda}{d\eta}(x)\Big)\,\lambda(dx)
    =
    \Ent(\lambda\mid\eta).
\end{equation*}

Now let \(\pi\) be a coupling of \(\widetilde\lambda\) and \(\mu\).
Then \(\pi\circ (T,T)^{-1}\) is a coupling of \(\lambda\) and \(\eta\), so by Lipschitz continuity of \(T\),
\begin{align*}
     & \int_{\R^d\times\R^d}
    |x-y|^p\,d(\pi\circ (T,T)^{-1})(x,y)          \\
     & \quad =
    \int_{\R^d\times\R^d}|T(u)-T(v)|^p\,d\pi(u,v) \\
     & \quad \le
    \Lip(T)^p
    \int |u-v|^p\,d\pi(u,v).
\end{align*}
Taking the infimum over all couplings \(\pi\) of \(\widetilde\lambda\) and \(\mu\) yields
\begin{equation*}
    W_p(\lambda,\eta)\le \Lip(T)\,W_p(\widetilde\lambda,\mu).
\end{equation*}
Since \(\mu\) satisfies \(T_p(C_T)\),
\begin{equation*}
    W_p(\widetilde\lambda,\mu)^2
    \le
    2C_T\,\Ent(\widetilde\lambda\mid\mu)
    =
    2C_T\,\Ent(\lambda\mid\eta).
\end{equation*}
Combining the two inequalities, we obtain
\begin{equation*}
    W_p(\lambda,\eta)^2
    \le
    \Lip(T)^2\,W_p(\widetilde\lambda,\mu)^2
    \le
    2C_T\,\Lip(T)^2\,\Ent(\lambda\mid\eta),
\end{equation*}
which proves that \(\eta\) satisfies \(T_p(C_T\,\Lip(T)^2)\).
\hfill\(\Box\)




















% ============================================================
% The averaged Jacobian
% ============================================================
\section{The averaged Jacobian}\label{sec:averaged-jacobian}

The preceding pathwise estimates use the quantity \(\mathcal K_t(\omega;K)\)
which controls the worst possible stretching of the flow on the compact set \(K\). Consider instead
the averaged Jacobian for the transported measure \(\mu_t=\mu_0\circ T_t^{-1}\),
\begin{equation*}
    \mathcal E_t := \int_{\R^d}\|J_t(u)\|_{\HS}^2\,\mu_0(du).
\end{equation*}
It measures average deformation of the flow over the initial mass, rather than
the maximal deformation on the whole support.


\begin{proposition}[Averaged Jacobian energy estimate]\label{prop:averaged-jacobian-energy}
    Assume that the flow is differentiable in the initial label and that there
    exist constants \(L_{\mathrm{sym}},\Gamma<\infty\) such that
    \begin{equation}\label{eq:one-sided-drift-jacobian}
        \langle z,\partial_xa(x,\mu)z\rangle
        \le
        L_{\mathrm{sym}}|z|^2,
        \qquad x,z\in\R^d,\ \mu\in\cP_2(\R^d),
    \end{equation}
    and
    \begin{equation}\label{eq:quadratic-diffusion-jacobian}
        \sum_{k=1}^d|\partial_x\sigma^{(k)}(x)z|^2
        \le
        \Gamma^2|z|^2,
        \qquad x,z\in\R^d.
    \end{equation}
    Then
    \begin{equation}\label{eq:averaged-jacobian-bound}
        \E\mathcal E_t
        \le
        d\exp\{(2L_{\mathrm{sym}}+\Gamma^2)t\}.
    \end{equation}
\end{proposition}

\noindent\emph{Proof.}
Recall that
\begin{equation*}
    dJ_t(u) = A_t(u)J_t(u)\,dt + \sum_{k=1}^d C_t^{(k)}(u)J_t(u)\,dB_t^k,
\end{equation*}
where
\begin{equation*}
    A_t(u)=\partial_xa(X_t(u),\mu_t), \qquad C_t^{(k)}(u)=\partial_x\sigma^{(k)}(X_t(u)).
\end{equation*}
Applying It\^o's formula to \(\|J_t(u)\|_{\HS}^2\) gives
\begin{align*}
    d\|J_t(u)\|_{\HS}^2 & =
    \bigg[2\langle J_t(u),A_t(u)J_t(u)\rangle_{\HS} +
    \sum_{k=1}^d\|C_t^{(k)}(u)J_t(u)\|_{\HS}^2 \bigg]dt                                                 \\
                        & \quad + 2\sum_{k=1}^d \langle J_t(u),C_t^{(k)}(u)J_t(u)\rangle_{\HS}\,dB_t^k.
\end{align*}
Since
\begin{equation*}
    \langle J,A_tJ\rangle_{\HS} = \sum_{i=1}^d\langle Je_i,A_tJe_i\rangle,
\end{equation*}
condition \cref{eq:one-sided-drift-jacobian} gives
\begin{equation*}
    2\langle J,A_tJ\rangle_{\HS} \le 2L_{\mathrm{sym}}\|J\|_{\HS}^2.
\end{equation*}
Similarly, \cref{eq:quadratic-diffusion-jacobian} gives
\begin{equation*}
    \sum_{k=1}^d\|C_t^{(k)}J\|_{\HS}^2
    =
    \sum_{i=1}^d\sum_{k=1}^d|C_t^{(k)}Je_i|^2
    \le
    \Gamma^2\|J\|_{\HS}^2.
\end{equation*}
Integrating with respect to \(\mu_0(du)\) and taking expectations removes the
martingale term, hence
\begin{equation*}
    \frac{d}{dt}\E\mathcal E_t \le (2L_{\mathrm{sym}}+\Gamma^2)\E\mathcal E_t.
\end{equation*}
Since \(J_0(u)=I\), we have \(\mathcal E_0=\|I\|_{\HS}^2=d\). Gronwall's
lemma proves \cref{eq:averaged-jacobian-bound}.
\hfill\(\Box\)


To estimate \(\E\sup_{u\in K}\|J_t(u)\|^p\), one must control increments \(J_t(u)-J_t(v)\), which requires assumptions on the spatial
regularity of \(\partial_xa\) and \(\partial\sigma\). The averaged quantity \(\mathcal E_t\) does not compare two labels. It only estimates each
\(J_t(u)\) and integrates in \(u\), so the one-sided bounds \cref{eq:one-sided-drift-jacobian,eq:quadratic-diffusion-jacobian}
are sufficient.





% ============================================================
% Internal fluctuations: pairwise energy and synchronization
% ============================================================
\section{Internal fluctuations}\label{sec:internal-fluctuations}

This section studies the pairwise energy
\begin{equation*}
    \mathcal D_t := \iint_{\R^d\times\R^d} |X_t(u)-X_t(v)|^2\,\mu_0(du)\mu_0(dv).
\end{equation*}

% ============================================================
% Internal fluctuations: pairwise energy and synchronization
% ============================================================
\subsection{Dissipativity assumptions}
Equivalently, if
\begin{equation*}
    \bar X_t:=\int_{\R^d}X_t(u)\,\mu_0(du),
\end{equation*}
and
\begin{equation*}
    \mathcal V_t := \int_{\R^d}|X_t(u)-\bar X_t|^2\,\mu_0(du),
\end{equation*}
then
\begin{equation*}
    \mathcal D_t=2\mathcal V_t.
\end{equation*}
Indeed,
\begin{align*}
    \mathcal D_t & = \iint |X_t(u)-X_t(v)|^2\,\mu_0(du)\mu_0(dv)                                                \\
                 & = \iint \big( |X_t(u)|^2+|X_t(v)|^2 -2\langle X_t(u),X_t(v)\rangle \big)\,\mu_0(du)\mu_0(dv) \\
                 & = 2\int |X_t(u)|^2\,\mu_0(du)- 2\left|\int X_t(u)\,\mu_0(du)\right|^2                        \\
                 & = 2\int |X_t(u)-\bar X_t|^2\,\mu_0(du)= 2\mathcal V_t.
\end{align*}
We assume the following dissipativity condition.
\begin{assumption}[Pairwise dissipativity]\label{ass:pairwise-dissipativity}
    There exists \(\lambda>0\) such that for every \(x,y\in\R^d\), every \(\theta\in\R^q\), and every
    \(\nu\in\cP_2(\R^d)\),
    \begin{align*}
         & 2\left\langle x-y, b(x,\theta)-b(y,\theta) + \int_{\R^d} \big[K(x,z,\theta)-K(y,z,\theta)\big] \nu(dz) \right\rangle \\
         & \qquad + \sum_{\ell=1}^m |\sigma_\ell(x,\theta)-\sigma_\ell(y,\theta)|^2 \le -\lambda |x-y|^2.
    \end{align*}
\end{assumption}

The dissipativity condition covers several basic models.

\begin{example}[Ornstein--Uhlenbeck SDEWI]
    Consider
    \begin{equation*}
        dX_t(u)= -\lambda_0\big(X_t(u)-m_t\big)\,dt + \sigma(X_t(u))\,dB_t, \qquad m_t=\int_{\R^d} X_t(v)\,\mu_0(dv).
    \end{equation*}
    If
    \begin{equation*}
        \|\sigma(x)-\sigma(y)\|_{\HS} \le L_\sigma |x-y|,
    \end{equation*}
    then
    \begin{equation*}
        \E\mathcal D_t \le e^{-(2\lambda_0-L_\sigma^2)t}\mathcal D_0.
    \end{equation*}
    Thus the cloud contracts whenever \(2\lambda_0>L_\sigma^2\).
\end{example}

\begin{example}[Gradient systems]
    Take
    \begin{equation*}
        b(x,\theta)=-\nabla V(x), \qquad K(x,y,\theta)=-\nabla W(x-y).
    \end{equation*}
    Then
    \begin{equation*}
        dX_t(u) =  \left[ -\nabla V(X_t(u)) - \int_{\R^d} \nabla W(X_t(u)-X_t(v))\,\mu_0(dv) \right]dt + \sigma(X_t(u))\,dB_t.
    \end{equation*}
    If
    \begin{equation*}
        D^2V\ge \rho I, \qquad D^2W\ge \kappa I, \qquad \|\sigma(x)-\sigma(y)\| \le L_\sigma |x-y|,
    \end{equation*}
    then
    \begin{equation*}
        \E\mathcal D_t \le e^{-(2\rho+2\kappa-L_\sigma^2)t}\mathcal D_0.
    \end{equation*}
\end{example}

\begin{example}[Nonlinear statistics]
    Let
    \begin{equation*}
        \Theta_t=\int_{\R^d} \psi(X_t(v))\,\mu_0(dv),
    \end{equation*}
    and consider
    \begin{equation*}
        dX_t(u) = -\lambda_0\big(X_t(u)-\Theta_t\big)\,dt + \sigma(X_t(u),\Theta_t)\,dB_t.
    \end{equation*}
    For two labels \(u,v\), the order parameter cancels in the drift
    difference:
    \begin{equation*}
        -\lambda_0\big(X_t(u)-\Theta_t\big) + \lambda_0\big(X_t(v)-\Theta_t\big) = -\lambda_0\big(X_t(u)-X_t(v)\big).
    \end{equation*}

\end{example}


\subsection{Proof}
\begin{theorem}[Synchronization under dissipativity assumptions]\label{thm:synchronization-dissipativity}
    Assume that \(\mu_0\in\cP_2(\R^d)\) and that \Cref{ass:pairwise-dissipativity} holds. Then
    \begin{equation*}
        \E\mathcal D_t \le e^{-\lambda t}\mathcal D_0.
    \end{equation*}

    In particular,
    \begin{equation*}
        \E W_2^2(\mu_t,\delta_{\bar X_t}) \le e^{-\lambda t}W_2^2(\mu_0,\delta_{\bar X_0}).
    \end{equation*}
    Thus the transported measure \(\mu_t\) concentrates around its random center of mass \(\bar X_t\).
\end{theorem}

% Proof idea:
% 1. Differentiate the stochastic flow.
\noindent\emph{Proof.}
% Proof roadmap: first prove pairwise contraction, then integrate over labels.
Fix two labels \(u,v\in\R^d\). Define the difference
\begin{equation*}
    Z_t^{u,v}:=X_t(u)-X_t(v).
\end{equation*}
We first prove that under \Cref{ass:pairwise-dissipativity}
\begin{equation*}
    \E|Z_t^{u,v}|^2 \le e^{-\lambda t}|u-v|^2,
\end{equation*}
then we integrate this estimate with respect to \(\mu_0(du)\mu_0(dv)\).

Writing \cref{eq:kernel-sdewi} for \(X_t(u)\) and for \(X_t(v)\) and then subtracting the two equations gives
\begin{equation*}
    dZ_t^{u,v} = G_t^{u,v}\,dt + \sum_{\ell=1}^m H_{\ell,t}^{u,v}\,dB_t^\ell,
\end{equation*}
where the drift coefficient
\begin{align*}
    G_t^{u,v} & := b(X_t(u),\Theta_t)-b(X_t(v),\Theta_t)                                                         \\
              & \quad+  \int_{\R^d} \big[ K(X_t(u),X_t(w),\Theta_t) - K(X_t(v),X_t(w),\Theta_t)\big]\,\mu_0(dw),
\end{align*}
and the diffusion coefficient is
\begin{equation*}
    H_{\ell,t}^{u,v} := \sigma_\ell(X_t(u),\Theta_t) - \sigma_\ell(X_t(v),\Theta_t).
\end{equation*}

Since \(\mu_t =\mu_0\circ X_t^{-1}\), we may rewrite the interaction term as
\begin{equation*}
    \int_{\R^d}K(X_t(u),X_t(w),\Theta_t)\,\mu_0(dw) = \int_{\R^d}K(X_t(u),z,\Theta_t)\,\mu_t(dz).
\end{equation*}
Thus
\begin{align*}
    G_t^{u,v} & := b(X_t(u),\Theta_t)-b(X_t(v),\Theta_t)                                               \\
              & \quad+ \int_{\R^d} \big[ K(X_t(u),z,\Theta_t) - K(X_t(v),z,\Theta_t) \big]\,\mu_t(dz).
\end{align*}

Applying It\^o's formula to \(|Z_t^{u,v}|^2\), we obtain
\begin{align*}
    d|Z_t^{u,v}|^2 & = 2\langle Z_t^{u,v},G_t^{u,v}\rangle\,dt + \sum_{\ell=1}^m |H_{\ell,t}^{u,v}|^2\,dt \\
                   & \quad+ 2\sum_{\ell=1}^m \langle Z_t^{u,v},H_{\ell,t}^{u,v}\rangle\,dB_t^\ell.
\end{align*}
Apply \Cref{ass:pairwise-dissipativity} with
\(x=X_t(u)\), \(y=X_t(v)\), \(\theta=\Theta_t\), and \(\nu=\mu_t\). Then
\begin{equation*} 2\langle Z_t^{u,v},G_t^{u,v}\rangle +  \sum_{\ell=1}^m |H_{\ell,t}^{u,v}|^2 \le -\lambda |Z_t^{u,v}|^2.
\end{equation*}
Therefore
\begin{equation}\label{eq:pairwise-contraction-differential}
    d|Z_t^{u,v}|^2 \le -\lambda |Z_t^{u,v}|^2\,dt + 2\sum_{\ell=1}^m \langle Z_t^{u,v},H_{\ell,t}^{u,v}\rangle\,dB_t^\ell.
\end{equation}

Set the new variable
\begin{equation*}
    Y_t^{u,v}:=e^{\lambda t}|Z_t^{u,v}|^2.
\end{equation*}
By It\^o's formula,
\begin{align*}
    dY_t^{u,v} & = e^{\lambda t} \left[ \lambda |Z_t^{u,v}|^2 + 2\langle Z_t^{u,v},G_t^{u,v}\rangle + \sum_{\ell=1}^m |H_{\ell,t}^{u,v}|^2 \right]dt \\
               & \quad+ 2e^{\lambda t} \sum_{\ell=1}^m \langle Z_t^{u,v},H_{\ell,t}^{u,v}\rangle\,dB_t^\ell.
\end{align*}
By the drift bound \cref{eq:pairwise-contraction-differential}, the drift term above is nonpositive. Hence
\begin{equation*}
    dY_t^{u,v} \le 2e^{\lambda t} \sum_{\ell=1}^m \langle Z_t^{u,v},H_{\ell,t}^{u,v}\rangle\,dB_t^\ell.
\end{equation*}
% Thus \(Y_t^{u,v}\) is a nonnegative local supermartingale.

To make this completely rigorous, introduce the localization sequence
\begin{equation*}
    \tau_R := \inf\{t\ge0: |X_t(u)|+|X_t(v)|+|\Theta_t|\ge R\}.
\end{equation*}
Since the solution is global and continuous, \(\tau_R\uparrow\infty\) almost surely as \(R\to\infty\). Stopping this inequality at \(\tau_R\), we get
\begin{equation*}
    Y_{t\wedge\tau_R}^{u,v} \le Y_0^{u,v} + 2\sum_{\ell=1}^m \int_0^{t\wedge\tau_R} e^{\lambda s} \langle Z_s^{u,v},H_{\ell,s}^{u,v}\rangle\,dB_s^\ell.
\end{equation*}
The stopped stochastic integral is a true martingale. Therefore, taking
expectations gives
\begin{equation*}
    \E Y_{t\wedge\tau_R}^{u,v} \le Y_0^{u,v}.
\end{equation*}
Since \(Y_0^{u,v}=|u-v|^2\), we have
\begin{equation*}
    \E\left[ e^{\lambda(t\wedge\tau_R)} |Z_{t\wedge\tau_R}^{u,v}|^2 \right] \le |u-v|^2.
\end{equation*}
Letting \(R\to\infty\) and using Fatou's lemma yields
\begin{equation*}
    \E\left[e^{\lambda t}|Z_t^{u,v}|^2\right] \le |u-v|^2.
\end{equation*}
Therefore
\begin{equation*}
    \E|X_t(u)-X_t(v)|^2 \le  e^{-\lambda t}|u-v|^2.
\end{equation*}
This proves the pairwise synchronization estimate.

Now integrate the pairwise synchronization estimate with respect to
\(\mu_0(du)\mu_0(dv)\). Since \(\mu_0\in\cP_2(\R^d)\),
\begin{equation*}
    \mathcal D_0 = \iint |u-v|^2\,\mu_0(du)\mu_0(dv)<\infty.
\end{equation*}
By Fubini's theorem,
\begin{align*}
    \E\mathcal D_t & = \E \iint |X_t(u)-X_t(v)|^2\,\mu_0(du)\mu_0(dv)      \\
                   & = \iint \E|X_t(u)-X_t(v)|^2\,\mu_0(du)\mu_0(dv)       \\
                   & \le e^{-\lambda t} \iint |u-v|^2\,\mu_0(du)\mu_0(dv).
\end{align*}
Hence \(\E\mathcal D_t\le e^{-\lambda t}\mathcal D_0\), proving
the first estimate in the theorem.

Therefore
\begin{equation*}
    \E\mathcal V_t = \frac12\E\mathcal D_t \le \frac12 e^{-\lambda t}\mathcal D_0 = e^{-\lambda t}\mathcal V_0.
\end{equation*}
This proves the internal fluctuation decay estimate.

Finally, since \(\mu_t =\mu_0\circ X_t^{-1}\),
\begin{equation*}
    \int_{\R^d}|x-\bar X_t|^2\,\mu_t(dx) = \int_{\R^d}|X_t(u)-\bar X_t|^2\,\mu_0(du) = \mathcal V_t.
\end{equation*}
For the Dirac measure \(\delta_{\bar X_t}\),
\begin{equation*}
    W_2^2(\mu_t,\delta_{\bar X_t}) = \int_{\R^d}|x-\bar X_t|^2\,\mu_t(dx).
\end{equation*}
Therefore
\begin{equation*}
    W_2^2(\mu_t,\delta_{\bar X_t})=\mathcal V_t.
\end{equation*}
Taking expectations and using the previous estimate gives
\begin{equation*}
    \E W_2^2(\mu_t,\delta_{\bar X_t}) \le e^{-\lambda t}W_2^2(\mu_0,\delta_{\bar X_0}).
\end{equation*}
\hfill\(\Box\)




\bibliographystyle{apalike}
\bibliography{reference}

\end{document}
